\documentclass[10pt,a4paper]{amsart}
\usepackage[margin=19mm]{geometry}
\usepackage{amsmath,amssymb,mathtools}
\usepackage[scr=rsfs,cal=euler]{mathalfa}
\usepackage{xfrac}
\usepackage[unicode,hidelinks,bookmarks=false]{hyperref}
\newcommand{\ie}{i.e.\ }
\newcommand{\cf}{cf.\ }
\newcommand{\resp}{respectively\ }
\newcommand{\Subsection}{\subsection}
\providecommand{\nobreakdash}{\nobreak\hbox{-}}
\setlength{\emergencystretch}{2em}
\multlinegap0pt
\usepackage[shortlabels]{enumitem}
\usepackage{amssymb}
\usepackage{mathtools}
\usepackage{tikz}
\newtheorem{lemma}{Lemma}
\newtheorem{theorem}{Theorem}
\newtheorem{proposition}{Proposition}
\newcommand{\AS}{\mathrm{AS}}
\newcommand{\diam}{\operatorname{diam}}
\title{Maximizing the Steklov eigenvalues on trees with a diameter constraint}
\author{Jiangdong Ai, Huiqiu Lin, Yongtang Shi}
\date{Source: arXiv:2604.12404v1 (CC BY 4.0)}
\begin{document}
\maketitle

\noindent\emph{Source and licence.} Maximizing the Steklov eigenvalues on trees with a diameter constraint. Version arXiv:2604.12404v1 (CC BY 4.0), distributed by arXiv under the Creative Commons Attribution 4.0 International licence, http://creativecommons.org/licenses/by/4.0/, as recorded in the arXiv licence metadata for this exact version and shown on the abstract page https://arxiv.org/abs/2604.12404v1. Full licence notice: \texttt{licenses/arxiv-2604.12404-CC-BY-4.0.txt}.\par\smallskip

\noindent\textit{Editorial scope.} Counted unit: the article's proposition prop:single-threshold (source0.tex lines 1516--1584). The article's complete original proof of it is delivered with it (source0.tex lines 1586--1778, one \texttt{proof} environment of 193 lines). No mathematics is written, repaired or re-derived: the statement and the proof are the article's own lines, in the article's own notation, and no retained line is substituted or re-flowed.  Retained uncounted as the source-local context the proof needs: lines 363--371; lines 1461--1493.  Nothing was omitted from any retained span, so the package carries no omission record.  No retained line contains a citation, so the wrapper carries no bibliography and no bibliography entry.  No retained line contains a figure environment, an image-inclusion command or drawing code, so no image is delivered and the package declares no asset dependency.  Editorial formatting only, in addition: an AMS \texttt{amsart} preamble that restates the article's own declarations and macros used by the retained lines, namely the environments \texttt{lemma}, \texttt{theorem}, \texttt{proposition} on separate counters, together with the standard packages those lines need.  The retained environments print in this document as Lemma 1, Theorem 1, Proposition 1 in order of appearance, whereas the article prints the same items with the numbering of its own full document.  The complete native source of the article as distributed in the arXiv e-print for this exact version is bundled unchanged as \texttt{sources/198399/source0.tex}.
\begin{lemma}\label{lem:spider-equation}
Let $T=S(\ell_1,\dots,\ell_b)$
be a spider with $\ell_1>\ell_2\ge \ell_3\ge \cdots \ge \ell_b\ge 1.$
Define $$
F_{\ell}(\lambda):=\sum_{i=1}^b \frac{1}{1-\ell_i\lambda}.
$$
Then the first nonzero Steklov eigenvalue $\lambda_2(T)$ is the unique root of $F_{\ell}$
in the interval $\left(\frac1{\ell_1},\frac1{\ell_2}\right).$
\end{lemma}

\begin{theorem}\label{thm:global-odd}
Let $D=2r+1\ge 3$ and let $n\ge D+1$.
Let $M=n-2r-2.$ Then the following hold.
\begin{enumerate}[label=\textnormal{(\roman*)}]
    \item If $M=0$, then the unique maximizer of $\lambda_2$ among all trees with order $n$ and
    diameter $D$ is the path $P_{D+1}$.
    \item Assume $M\ge 1$, and define
    $$
    s=\left\lceil \frac{r}{2}\right\rceil,
    \qquad
    q_-=\max\!\left\{1,\left\lfloor \frac{M}{s}\right\rfloor\right\},
    \qquad
    q_+=\left\lceil \frac{M}{s}\right\rceil.
    $$
    Write
    $$
    M=q_-c_-+t_-,\qquad 0\le t_-<q_-,
    $$
    $$
    M=q_+c_++t_+,\qquad 0\le t_+<q_+.
    $$
    Then
    \begin{equation}\label{eq:global-two-candidates}
    \max_{\substack{|V(T)|=n\\ \diam(T)=D}} \lambda_2(T)
    =
    \max\Bigl\{\lambda_2\bigl(\AS(r,q_-+2,c_-,t_-)\bigr),\,
    \lambda_2\bigl(\AS(r,q_++2,c_+,t_+)\bigr)\Bigr\}.
    \end{equation}
    Every maximizer is isomorphic either to $\AS(r,q_-+2,c_-,t_-)$ or to $\AS(r,q_++2,c_+,t_+)$.
    In particular, if the two values on the right-hand side of \eqref{eq:global-two-candidates}
    are different, then the maximizer is unique.
\end{enumerate}
\end{theorem}

\begin{proposition}\label{prop:single-threshold}
Fix $r\ge 2$ and let $s=\lceil r/2\rceil$.
Let $1\le t\le s-1$, let $k\ge s$, and let $M=ks+t$.
Then
$$
q_-=\left\lfloor \frac{M}{s}\right\rfloor=k,
\qquad
q_+=\left\lceil \frac{M}{s}\right\rceil=k+1.
$$
Hence the two candidates from Theorem~\ref{thm:global-odd} are
$$
A_{k,t}=\AS(r,k+2,s,t)
=S\bigl(r+1,r,\underbrace{s+1,\dots,s+1}_{t},\underbrace{s,\dots,s}_{k-t}\bigr),
$$
and
$$
B_{k,t}=\AS(r,k+3,s-1,k+t-s+1)
=S\bigl(r+1,r,\underbrace{s,\dots,s}_{k+t-s+1},\underbrace{s-1,\dots,s-1}_{s-t}\bigr).
$$
Let $F_{A_{k,t}}$ and $F_{B_{k,t}}$ be the corresponding root functions from Lemma~\ref{lem:spider-equation}.
Then, on $I_r=\left(\frac1{r+1},\frac1r\right)$,
we have
\begin{equation}\label{eq:single-threshold-difference}
F_{B_{k,t}}(\lambda)-F_{A_{k,t}}(\lambda)
=
\frac{P_{s,t}(\lambda)}{(1-(s-1)\lambda)(1-s\lambda)(1-(s+1)\lambda)},
\end{equation}
where
$P_{s,t}(\lambda)=1-3s\lambda+(2s^2+s-2t-1)\lambda^2.$
Moreover, the following hold.
\begin{enumerate}[label=\textnormal{(\roman*)}]
    \item If $r=2s$ and $2t\le s-1$, then
    $$
    \lambda_2(A_{k,t})>\lambda_2(B_{k,t})
    \qquad\text{for every }k\ge s.
    $$
    \item If $r=2s-1$ and $2t\ge s-1$, then
    $$
    \lambda_2(B_{k,t})>\lambda_2(A_{k,t})
    \qquad\text{for every }k\ge s.
    $$
    \item In the remaining cases, namely
    $$
    r=2s,\ 2t\ge s,
    \qquad\text{or}\qquad
    r=2s-1,\ 2t\le s-2,
    $$
    the polynomial $P_{s,t}$ has a unique root $\zeta_{s,t}$ in $I_r$.
    Define
    $$
    \kappa_{r,t}
    =-(1-s\zeta_{s,t})\left(
    \frac{1}{1-(r+1)\zeta_{s,t}}
    +\frac{1}{1-r\zeta_{s,t}}
    +\frac{t-s+1}{1-s\zeta_{s,t}}
    +\frac{s-t}{1-(s-1)\zeta_{s,t}}
    \right).
    $$
    Then, for every $k\ge s$,
    $$
    \lambda_2(A_{k,t})>\lambda_2(B_{k,t})\iff k>\kappa_{r,t},
    $$
    $$
    \lambda_2(B_{k,t})>\lambda_2(A_{k,t})\iff k<\kappa_{r,t},
    $$
    and equality holds if and only if $k=\kappa_{r,t}\in\mathbb Z$.
    In particular, for each fixed residue class $t$, there is at most one value of $k$ for which the two candidates tie.
\end{enumerate}
\end{proposition}

\begin{proof}
The identities $q_-=k$ and $q_+=k+1$ are immediate from $M=ks+t$ with $1\le t\le s-1$.
The displayed forms of $A_{k,t}$ and $B_{k,t}$ follow directly from the definition of the generalized almost seesaw trees.

Write
$$
\alpha_{k,t}=\lambda_2(A_{k,t}),
\qquad
\beta_{k,t}=\lambda_2(B_{k,t}).
$$
By Lemma~\ref{lem:spider-equation}, both $F_{A_{k,t}}$ and $F_{B_{k,t}}$ are strictly increasing on $I_r$, and $\alpha_{k,t},\beta_{k,t}$ are their unique roots there.

A direct computation gives
\begin{align*}
F_{A_{k,t}}(\lambda)
&=
\frac{1}{1-(r+1)\lambda}
+\frac{1}{1-r\lambda}
+\frac{t}{1-(s+1)\lambda}
+\frac{k-t}{1-s\lambda},\\
F_{B_{k,t}}(\lambda)
&=
\frac{1}{1-(r+1)\lambda}
+\frac{1}{1-r\lambda}
+\frac{k+t-s+1}{1-s\lambda}
+\frac{s-t}{1-(s-1)\lambda}.
\end{align*}
Subtracting the two expressions yields \eqref{eq:single-threshold-difference}.

Since $s\le r$, all three denominator factors in \eqref{eq:single-threshold-difference} are positive on $I_r$.
Hence the sign of $F_{B_{k,t}}-F_{A_{k,t}}$ on $I_r$ is exactly the sign of $P_{s,t}$.

We first show that $P_{s,t}$ is strictly decreasing on $I_r$.
Indeed,
$$
P'_{s,t}(\lambda)=-3s+2(2s^2+s-2t-1)\lambda.
$$
If $r=2s$, then $\lambda<1/r=1/(2s)$, and therefore
$$
P'_{s,t}(\lambda)
<
-3s+\frac{2(2s^2+s-2t-1)}{2s}
=
-s+1-\frac{2t+1}{s}
<0.
$$
If $r=2s-1$, then $\lambda<1/r=1/(2s-1)$, and hence
$$
P'_{s,t}(\lambda)
<
-3s+\frac{2(2s^2+s-2t-1)}{2s-1}
=
\frac{-2s^2+5s-4t-2}{2s-1}
<0.
$$
Thus $P_{s,t}$ is strictly decreasing on $I_r$ in all cases.

We now evaluate $P_{s,t}$ at the endpoints of $I_r$.

\medskip
\noindent
{\it Case 1: $r=2s$.}
Then
$$
P_{s,t}\!\left(\frac1{2s+1}\right)=\frac{2(s-t)}{(2s+1)^2}>0,
$$
and
$$
P_{s,t}\!\left(\frac1{2s}\right)=\frac{s-2t-1}{4s^2}.
$$
If $2t\le s-1$, then both endpoint values are nonnegative, and strict decrease implies
$$
P_{s,t}(\lambda)>0\qquad \text{for all }\lambda\in I_r.
$$
Hence $F_{B_{k,t}}(\lambda)-F_{A_{k,t}}(\lambda)>0$ on $I_r$.
Evaluating at $\lambda=\beta_{k,t}$ gives
$$
0-F_{A_{k,t}}(\beta_{k,t})>0,
$$
so $F_{A_{k,t}}(\beta_{k,t})<0$.
Since $F_{A_{k,t}}$ is strictly increasing and vanishes at $\alpha_{k,t}$, we obtain $\beta_{k,t}<\alpha_{k,t}$, that is,
$$
\lambda_2(A_{k,t})>\lambda_2(B_{k,t}).
$$
This proves \textnormal{(i)}.

If instead $2t\ge s$, then
$$
P_{s,t}\!\left(\frac1{2s+1}\right)>0,
\qquad
P_{s,t}\!\left(\frac1{2s}\right)<0.
$$
Since $P_{s,t}$ is strictly decreasing, it has a unique root $\zeta_{s,t}$ in $I_r$.

\medskip
\noindent
{\it Case 2: $r=2s-1$.}
Then
$$
P_{s,t}\!\left(\frac1{2s}\right)=\frac{s-2t-1}{4s^2},
$$
while
$$
P_{s,t}\!\left(\frac1{2s-1}\right)= -\frac{2t}{(2s-1)^2}<0.
$$
If $2t\ge s-1$, then both endpoint values are nonpositive, and strict decrease implies
$$
P_{s,t}(\lambda)<0\qquad \text{for all }\lambda\in I_r.
$$
Hence $F_{B_{k,t}}(\lambda)-F_{A_{k,t}}(\lambda)<0$ on $I_r$.
Evaluating at $\lambda=\alpha_{k,t}$ gives
$$
F_{B_{k,t}}(\alpha_{k,t})-0<0,
$$
so $F_{B_{k,t}}(\alpha_{k,t})<0$.
Since $F_{B_{k,t}}$ is strictly increasing and vanishes at $\beta_{k,t}$, we get $\alpha_{k,t}<\beta_{k,t}$, that is,
$$
\lambda_2(B_{k,t})>\lambda_2(A_{k,t}).
$$
This proves \textnormal{(ii)}.

If instead $2t\le s-2$, then
$$
P_{s,t}\!\left(\frac1{2s}\right)>0,
\qquad
P_{s,t}\!\left(\frac1{2s-1}\right)<0,
$$
so again $P_{s,t}$ has a unique root $\zeta_{s,t}$ in $I_r$.

\medskip
We are thus left with the two threshold cases, in which $P_{s,t}$ has a unique root $\zeta_{s,t}\in I_r$.
Since $P_{s,t}$ is strictly decreasing, we have
$$
P_{s,t}(\lambda)>0\iff \lambda<\zeta_{s,t},
\qquad
P_{s,t}(\lambda)<0\iff \lambda>\zeta_{s,t}.
$$
Equivalently,
$$
F_{B_{k,t}}(\lambda)-F_{A_{k,t}}(\lambda)>0\iff \lambda<\zeta_{s,t}.
$$

Now
$$
\lambda_2(A_{k,t})>\lambda_2(B_{k,t})
\iff \alpha_{k,t}>\beta_{k,t}
\iff F_{A_{k,t}}(\beta_{k,t})<0
\iff F_{B_{k,t}}(\beta_{k,t})-F_{A_{k,t}}(\beta_{k,t})>0,
$$
and hence
$$
\lambda_2(A_{k,t})>\lambda_2(B_{k,t})\iff \beta_{k,t}<\zeta_{s,t}.
$$
Since $F_{B_{k,t}}$ is strictly increasing and vanishes at $\beta_{k,t}$, this is equivalent to
$$
F_{B_{k,t}}(\zeta_{s,t})>0.
$$
Likewise,
$$
\lambda_2(B_{k,t})>\lambda_2(A_{k,t})\iff F_{B_{k,t}}(\zeta_{s,t})<0,
$$
and equality holds if and only if $F_{B_{k,t}}(\zeta_{s,t})=0$.

Finally, we compute $F_{B_{k,t}}(\zeta_{s,t})$:
\begin{align*}
F_{B_{k,t}}(\zeta_{s,t})
&=
\frac{1}{1-(r+1)\zeta_{s,t}}
+\frac{1}{1-r\zeta_{s,t}}
+\frac{k+t-s+1}{1-s\zeta_{s,t}}
+\frac{s-t}{1-(s-1)\zeta_{s,t}}\\
&=
\frac{k}{1-s\zeta_{s,t}}
+\left(
\frac{1}{1-(r+1)\zeta_{s,t}}
+\frac{1}{1-r\zeta_{s,t}}
+\frac{t-s+1}{1-s\zeta_{s,t}}
+\frac{s-t}{1-(s-1)\zeta_{s,t}}
\right)\\
&=
\frac{k-\kappa_{r,t}}{1-s\zeta_{s,t}}.
\end{align*}
Because $1-s\zeta_{s,t}>0$, the sign of $F_{B_{k,t}}(\zeta_{s,t})$ is exactly the sign of $k-\kappa_{r,t}$.
Therefore
$$
\lambda_2(A_{k,t})>\lambda_2(B_{k,t})\iff k>\kappa_{r,t},
$$
$$
\lambda_2(B_{k,t})>\lambda_2(A_{k,t})\iff k<\kappa_{r,t},
$$
and equality holds if and only if $k=\kappa_{r,t}\in\mathbb Z$.
This proves \textnormal{(iii)}.
\end{proof}

\end{document}
